\documentclass[11pt,respuestas,a4]{aleph-examen}
% \documentclass[11pt,a4]{aleph-examen}

% -- Paquetes extra
\usepackage{aleph-comandos}
\usepackage{systeme}

% -- Datos
\institucion{Facultad de Ciencias Exactas, Naturales y Ambientales}
\carrera{Catálogo STEM}
\asignatura{Matemática III}
\tema{Repaso -- Criterio 2.2}
\autor{Andrés Merino}
\fecha{Periodo 2026-1}

\logouno[0.16\textwidth]{Logos/logoPUCE_04_ac}
\definecolor{colortext}{HTML}{1957d1}
\definecolor{colordef}{HTML}{1957d1}
\fuente{montserrat}


\begin{document}

\encabezado


\begin{preguntas}

%%%%%%%%%%%%%%%%%%%%%%%%%%%%%%%%%%%%%%%%%%
\item
Considere la función
\[
    f(x,y) = \begin{cases}
        \dfrac{x^2 y}{x^4 + y^2} & \text{si } (x,y)\neq(0,0),\\[6pt]
        0 & \text{si } (x,y)=(0,0).
    \end{cases}
\]
\begin{enumerate}
    \item Calcule los límites iterados cuando $(x,y)\to(0,0)$.
    \item Calcule el límite cuando $(x,y)\to(0,0)$ por la trayectoria $y=x^2$.
    \item Calcule el límite por la trayectoria $y=0$.
    \item ¿Es $f$ continua en el origen? Justifique a partir de los incisos anteriores.
\end{enumerate}

\begin{respuesta}
\begin{enumerate}
    \item Calculemos los límites iterados. Primero cuando \(y\to0\) y luego cuando \(x\to0\):
    \[
        \lim_{x\to 0}\left(\lim_{y\to 0} f(x,y)\right)
        =
        \lim_{x\to 0}\left(\lim_{y\to 0} \frac{x^2y}{x^4+y^2}\right)
        =
        \lim_{x\to 0}\left(\frac{x^2(0)}{x^4+(0)^2}\right)
        =
        \lim_{x\to 0} 0
        =
        0.
    \]

    Ahora cuando \(x\to0\) y luego cuando \(y\to0\):
    \[
        \lim_{y\to 0}\left(\lim_{x\to 0} f(x,y)\right)
        =
        \lim_{y\to 0}\left(\lim_{x\to 0} \frac{x^2y}{x^4+y^2}\right)
        =
        \lim_{y\to 0}\left(\frac{(0)^2y}{(0)^4+y^2}\right)
        =
        \lim_{y\to 0} 0
        =
        0.
    \]

    \item Consideremos la trayectoria \(y=x^2\). Entonces:
    \[
        f(x,x^2)
        =
        \frac{x^2(x^2)}{x^4+(x^2)^2}
        =
        \frac{x^4}{x^4+x^4}
        =
        \frac{x^4}{2x^4}
        =
        \frac{1}{2}.
    \]
    Por tanto,
    \[
        \lim_{x\to 0} f(x,x^2)
        =
        \frac{1}{2}.
    \]

    \item Consideremos la trayectoria \(y=0\). Entonces:
    \[
        f(x,0)
        =
        \frac{x^2(0)}{x^4+0^2}
        =
        0.
    \]
    Por tanto,
    \[
        \lim_{x\to 0} f(x,0)
        =
        0.
    \]

    \item  Como los límites por distintas trayectorias son diferentes, se concluye que:
    \[
        \lim_{(x,y)\to(0,0)} f(x,y)
    \]
    no existe. Por lo tanto, \(f\) no es continua en \((0,0)\).\qedhere
\end{enumerate}
\end{respuesta}

\item
Dada $f(x,y,z) = x^2 \ln(y+z)$ y el punto $P=(2, 1, 0)$:
\begin{enumerate}
    \item Calcule $\nabla f(P)$, interprete el resultado.
    \item Calcule la derivada direccional de $f$ en $P$ en la dirección del vector $\overrightarrow{PQ}$, donde $Q=(4,2,1)$.
    \item ¿En qué dirección la tasa de cambio de $f$ en $P$ es igual a cero? ¿Qué representa geométricamente?
\end{enumerate}

\begin{respuesta}
\begin{enumerate}
    \item Calculemos el gradiente de \(f\). Para ello, usamos las derivadas parciales:
    \[
        \frac{\partial f}{\partial x}(x,y,z)
        =
        2x\ln(y+z),
        \qquad
        \frac{\partial f}{\partial y}(x,y,z)
        =
        x^2\frac{1}{y+z},
        \qquad
        \frac{\partial f}{\partial z}(x,y,z)
        =
        x^2\frac{1}{y+z}.
    \]
    Por lo tanto:
    \[
        \nabla f(x,y,z)
        =
        \left(
            2x\ln(y+z),\ 
            \frac{x^2}{y+z},\ 
            \frac{x^2}{y+z}
        \right).
    \]
    Evaluamos en \(P=(2,1,0)\):
    \[
        \nabla f(P)
        =
        \left(
            2(2)\ln(1+0),\ 
            \frac{2^2}{1+0},\ 
            \frac{2^2}{1+0}
        \right)
        =
        (0,4,4).
    \]
    El gradiente \(\nabla f(P)=(0,4,4)\) indica la dirección de máximo crecimiento de \(f\) en el punto \(P\).

    \item Primero calculamos el vector \(\overrightarrow{PQ}\):
    \[
        u = \overrightarrow{PQ}
        =
        Q-P
        =
        (4,2,1)-(2,1,0)
        =
        (2,1,1).
    \]
    La derivada direccional de \(f\) en \(P\) en la dirección de \(u\) es:
    \[
        f'(P; u)
        =
        \nabla f(P)\cdot u.
    \]
    Por tanto:
    \[
        f'(P; u)
        =
        (0,4,4)\cdot(2,1,1)
        =
        8.
    \]

    \item La tasa de cambio de \(f\) en \(P\) en la dirección de un vector unitario \(u=(a,b,c)\) es:
    \[
        f'(P; u)
        =
        \nabla f(P)\cdot u.
    \]
    Queremos que esta tasa de cambio sea igual a cero:
    \[
        \nabla f(P)\cdot u
        =
        0.
    \]
    Como \(\nabla f(P)=(0,4,4)\), entonces:
    \[
        (0,4,4)\cdot(a,b,c)
        =
        0.
    \]
    De donde:
    \[
        4b+4c
        =
        0.
    \]
    Por lo tanto:
    \[
        b+c=0.
    \]
    Así, la tasa de cambio es cero en cualquier dirección unitaria \((a,b,c)\) que satisfaga esta condición, por ejemplo, una dirección posible es:
    \[
        (1,0,0).
    \]
    Geométricamente, estas direcciones son perpendiculares al gradiente \(\nabla f(P)\). Por tanto, son direcciones tangentes a la superficie de nivel de \(f\) que pasa por \(P\), es decir, direcciones en las que \(f\) no cambia instantáneamente.\qedhere
\end{enumerate}
\end{respuesta}

\item
La ecuación $x^2+2y^2z-z^3-2=0$ define implícitamente \(z\) en función de \((x,y)\) cerca del punto \((1,1,1)\).
\begin{enumerate}
    \item Calcule \(\dfrac{\partial z}{\partial x}\) y \(\dfrac{\partial z}{\partial y}\) en \((1,1,1)\) usando el teorema de la función implícita.
\end{enumerate}

\begin{respuesta}
Definimos
    \[
        F(x,y,z)
        =
        x^2+2y^2z-z^3-2.
    \]
    Calculamos las derivadas parciales:
    \[
        \frac{\partial F}{\partial x}(x,y,z)
        =
        2x,
        \qquad
        \frac{\partial F}{\partial y}(x,y,z)
        =
        4yz,
        \qquad
        \frac{\partial F}{\partial z}(x,y,z)
        =
        2y^2-3z^2.
    \]
    Evaluando en \((1,1,1)\):
    \[
        \frac{\partial F}{\partial x}(1,1,1)
        =
        2,
        \qquad
        \frac{\partial F}{\partial y}(1,1,1)
        =
        4,
        \qquad
        \frac{\partial F}{\partial z}(1,1,1)
        =
        -1.
    \]
    Como
    \[
        \frac{\partial F}{\partial z}(1,1,1)\neq 0,
    \]
    por el teorema de la función implícita se tendría
    \[
        \frac{\partial z}{\partial x}
        =
        -\frac{\dfrac{\partial F}{\partial x}}{\dfrac{\partial F}{\partial z}},
        \qquad
        \frac{\partial z}{\partial y}
        =
        -\frac{\dfrac{\partial F}{\partial y}}{\dfrac{\partial F}{\partial z}}.
    \]
    Por tanto,
    \[
        \frac{\partial z}{\partial x}(1,1)
        =
        -\frac{2}{-1}
        =
        2,\qquad\text{y}\qquad
        \frac{\partial z}{\partial y}(1,1)
        =
        -\frac{4}{-1}
        =
        4.\qedhere
    \]
\end{respuesta}


\item
Sea $f(x,y) = e^{x}\cos(y)$.
\begin{enumerate}
    \item Calcule la matriz Hessiana $H_f(0,0)$.
    \item Escriba el polinomio de Taylor de orden 2 de $f$ alrededor de $(0,0)$ (aproximación cuadrática).
    \item Use el polinomio obtenido para estimar $f(0.1\ ,\ 0.1)$ y compárelo con el valor exacto. ¿Es buena la aproximación?
\end{enumerate}

\begin{respuesta}
\begin{enumerate}
    \item Calculemos las derivadas parciales de segundo orden. Primero:
    \[
        \frac{\partial f}{\partial x}(x,y)
        =
        e^x\cos(y),
        \qquad
        \frac{\partial f}{\partial y}(x,y)
        =
        -e^x\sen(y).
    \]
    Luego:
    \[
        \frac{\partial^2 f}{\partial x^2}(x,y)
        =
        e^x\cos(y),
        \qquad
        \frac{\partial^2 f}{\partial x\partial y}(x,y)
        =
        -e^x\sen(y),
        \qquad
        \frac{\partial^2 f}{\partial y^2}(x,y)
        =
        -e^x\cos(y).
    \]
    Evaluamos en \((0,0)\):
    \[
        \frac{\partial^2 f}{\partial x^2}(0,0)
        =
        1,
        \qquad
        \frac{\partial^2 f}{\partial x\partial y}(0,0)
        =
        0,
        \qquad
        \frac{\partial^2 f}{\partial y^2}(0,0)
        =
        -1.
    \]
    Por lo tanto:
    \[
        H_f(0,0)
        =
        \begin{pmatrix}
            1 & 0\\
            0 & -1
        \end{pmatrix}.
    \]

    \item Usemos la expresión:
    \[
        F(x,y)
        =
        f(a)
        +
        \nabla f(a)\cdot ((x,y)-a)
        +
        \frac{1}{2!}[(x,y)-a]^T H_f(a)[(x,y)-a].
    \]
    En este caso:
    \[
        a=(0,0),
        \qquad
        (x,y)-a=(x,y).
    \]
    Además:
    \[
        f(0,0)
        =
        1,
        \qquad
        \nabla f(0,0)
        =
        (1,0),
        \qquad
        H_f(0,0)
        =
        \begin{pmatrix}
            1 & 0\\
            0 & -1
        \end{pmatrix}.
    \]
    Entonces:
    \[
        F(x,y)
        =
        1
        +
        (1,0)\cdot(x,y)
        +
        \frac{1}{2}
        \begin{pmatrix}
            x & y
        \end{pmatrix}
        \begin{pmatrix}
            1 & 0\\
            0 & -1
        \end{pmatrix}
        \begin{pmatrix}
            x\\
            y
        \end{pmatrix}.
    \]
    Por lo tanto, la aproximación cuadrática es
    \[
        F(x,y)
        =
        1+x+\frac{1}{2}(x^2-y^2).
    \]

    \item Usamos el polinomio obtenido para estimar \(f(0.1,\ 0.1)\):
    \[
        F(0.1,\ 0.1)
        =
        1+0.1+\frac{1}{2}\left((0.1)^2-(0.1)^2\right)
        =
        1.1.
    \]
    El valor exacto es:
    \[
        f(0.1,\ 0.1)
        =
        e^{0.1}\cos(0.1)
        \approx
        1.09965.
    \]
    Por tanto, el error absoluto es:
    \[
        |f(0.1,\ 0.1)-F(0.1,\ 0.1)|
        \approx
        |1.09965-1.1|
        =
        0.00035.
    \]
    La aproximación es buena, pues el error es pequeño.\qedhere
\end{enumerate}
\end{respuesta}

\item
Dado el campo vectorial $F(x,y) = (x^2y,\; xy^2)$:
\begin{enumerate}
    \item Calcule $\mathrm{div}(F)$ y $\mathrm{rot}(F)$.
    \item Determine las regiones del espacio donde $F$ tiene fuentes y donde tiene sumideros.
    \item Determine las regiones donde $\mathrm{rot}(F)=0$.
\end{enumerate}

\begin{respuesta}
\begin{enumerate}
    \item Sea \(F(x,y)=(f_1(x,y),f_2(x,y))\), donde:
    \[
        f_1(x,y)=x^2y,
        \qquad
        f_2(x,y)=xy^2.
    \]
    La divergencia de \(F\) es:
    \[
        \mathrm{div}(F)(x,y)
        =
        \frac{\partial f_1}{\partial x}(x,y)
        +
        \frac{\partial f_2}{\partial y}(x,y).
    \]
    Entonces:
    \[
        \mathrm{div}(F)(x,y)
        =
        2xy+2xy
        =
        4xy.
    \]

    El rotacional de \(F\) en el plano es:
    \[
        \mathrm{rot}(F)(x,y)
        =
        \frac{\partial f_2}{\partial x}(x,y)
        -
        \frac{\partial f_1}{\partial y}(x,y).
    \]
    Por tanto:
    \[
        \mathrm{rot}(F)(x,y)
        =
        y^2-x^2.
    \]

    \item Para determinar dónde hay fuentes o sumideros, analizamos el signo de la divergencia.
    Hay fuentes donde:
    \[
        \mathrm{div}(F)(x,y)>0.
    \]
    Esto ocurre cuando:
    \[
        xy>0.
    \]
    Es decir, en los cuadrantes I y III. Hay sumideros donde:
    \[
        \mathrm{div}(F)(x,y)<0.
    \]
    Esto ocurre cuando:
    \[
        xy<0.
    \]
    Es decir, en los cuadrantes II y IV. Además, sobre los ejes:
    \[
        xy=0,
    \]
    por lo que:
    \[
        \mathrm{div}(F)(x,y)=0.
    \]
    En los ejes no hay fuentes ni sumideros locales.

    \item Para determinar dónde \(\mathrm{rot}(F)=0\), resolvemos:
    \[
        y^2-x^2=0.
    \]
    Entonces:
    \[
        (y-x)(y+x)=0.
    \]
    Por lo tanto:
    \[
        y=x
        \qquad
        \text{o}
        \qquad
        y=-x.
    \]

    Así, el rotacional se anula sobre las rectas:
    \[
        y=x
        \qquad
        \text{y}
        \qquad
        y=-x.
    \]
    En estas regiones, el campo no presenta rotación instantánea local.\qedhere
\end{enumerate}
\end{respuesta}

\end{preguntas}

\end{document}
